\documentclass[11pt]{article}
\usepackage{amsmath,amssymb,amsthm,mathtools,hyperref}

\begin{document}
\section*{General disease ecology}
\label{prob:200}
\label{rank:242}
\par\smallskip

\noindent\textbf{Domain.} Ecological epidemiology; host-community disease. \quad \textbf{Retained status/source context.} Open cross-system calibration and transfer problem. Multi-host next-generation matrices and many host--vector--environment models are established; their existence is not unresolved.
\paragraph{Definitions and assumptions.} Choose a specified host community with susceptible/infected classes, optional vector classes and an environmental infectious-stage compartment. Abundances obey a finite-parameter ecological model including resources and interspecies competition. At a constant disease-free equilibrium, matrix $F$ records new infection rates and $V$ transitions/removal of infected states; assumptions make $V$ a nonsingular M-matrix and $FV^{-1}$ nonnegative. For each instance, declare the ordered infected-state vector $z$, including the vector and environmental infectious states used, and the decomposition of its linearized dynamics as $\dot z=(F-V)z$. Specify the evaluation equilibrium and its selection if more than one disease-free equilibrium is available, and state the ecological and epidemiological hypotheses supporting the intended invasion interpretation. The allocation of environmental production between new-infection and transition terms, and hence the biological generation convention, is not selected here and remains to be declared; matrix positivity alone does not select it. No nonlinear invasion conclusion at $R_0=1$ is supplied by the displayed spectral-radius formula alone. Contact mode, shedding, environmental persistence and detectability are measured separately. Define infection burden $B_T$ as cumulative infected host-time over horizon $T$. A concrete instance must declare a common burden readout under both intervention laws: the contributing host population and infected classes, the conversion from ecological state variables to infected host abundance over a declared spatial extent, and the time origin and units for integration over $[0,T]$. The choice between total community burden and burden in a fixed focal host population remains unspecified; neither is selected by this formulation. The disease-free ecological attractor is estimated from the same community assembly considered for invasion. Adding a host changes both transmission opportunities and ecological densities, so holding all background abundances fixed represents a separate intervention from allowing the ecosystem to re-equilibrate.
\paragraph{Problem statement.} Predict how a declared host-community intervention changes invasion and burden:
\[
 R_0=\rho(FV^{-1}),\qquad
 \Delta B_T=\mathbb{E}^{\mathcal I}[B_T]-\mathbb{E}^{\mathcal I_0}[B_T],
\]
with an uncertainty interval narrow enough to distinguish the sign of $\Delta B_T$ where identifiable. Validate transport to a new community with measured ecological covariates, without substituting richness alone for host identity or competence.
\paragraph{Additional binding assumptions and resolution criterion.} Interpret $\mathcal I$ as a fully specified intervention-generated law, not as ordinary conditioning on the observed presence of a host assemblage. The next-generation spectral-radius threshold is an established result under its assumptions, whereas burden is a nonlinear finite-horizon quantity requiring initial infection and observation laws. A sign-identification target should permit an inconclusive result when zero remains in the joint uncertainty interval.

\paragraph{Alternative formulations.} (i) Under a constant disease-free equilibrium estimate $R_0$ and compare density- with frequency-dependent transmission. (ii) Under seasonal environments or ecological feedback use a periodic invasion exponent and finite-population outbreak/extinction probabilities; $R_0>1$ does not itself determine endemic burden. A periodic invasion exponent requires a declared periodic disease-free trajectory and period; ecological feedback alone does not imply periodicity. This periodic formulation therefore does not specify a nonperiodic invasion target. Outbreak and extinction probabilities require separately declared events, initial laws and time horizons, as applicable under the binding catalogue conventions.
\paragraph{Mathematical-status qualification.} The restricted calculation, identity, obstruction or construction identified in this entry is not itself a new open conjecture. The residual target is the explicitly delimited extension, identification, implementation or independent biological validation; an already known special case does not satisfy that residual target.

\paragraph{Obstructions and partial progress.} The ecological next-generation framework already handles interacting host/non-host species. A 22-year experiment documents a reversal in the diversity--disease relationship, so a universal dilution sign is untenable. Hidden reservoirs, contact heterogeneity and environmental stages can produce similar observed prevalence. The unresolved target is identifiable multi-route transmission and community-intervention prediction with ecological feedback, rather than a single general formula covering every ecosystem.
\paragraph{Sources (supplied provenance; not all rechecked).}
\begin{itemize}
\item Dobson, Holt and Tilman (eds.), Unsolved Problems in Ecology (2020). \url{https://www.jstor.org/stable/j.ctvs9fh2n}. Supplied book; contents and chapter previews checked, full text unavailable. Historical agenda, not present-day certification. Relevant listed chapter: Michael E. Hochberg, Six Wedges to Curing Disease.
\item Roberts and Heesterbeek, Characterizing the next-generation matrix and basic reproduction number in ecological epidemiology (2013). \url{https://pubmed.ncbi.nlm.nih.gov/23086599/}. Established ecological next-generation framework and specific examples.
\item The effect of diversity on disease reverses from dilution to amplification in a 22-year biodiversity x N x CO2 experiment (2024). \url{https://www.nature.com/articles/s41598-024-60725-z}. Primary experimental counterexample to a universal diversity sign.
\end{itemize}


\section*{Binding catalogue conventions}
This edition revises both supplied catalogues and preserves original problem numbers 1--330. It adds one globally unique label \verb|\label{rank:x}| per problem, where smaller $x$ means greater provisional difficulty. A rank label is an anchor at that problem's section: \verb|\ref{rank:x}| returns the section number, not the literal difficulty number. The visible \texttt{rank:x} text and the index give the difficulty number itself.

The entries are \emph{parameterized mathematical research schemes}, including empirical prediction, causal identification, historical reconstruction, theory and engineering. They are not 330 certified unsolved mathematical conjectures. A rigorous instance fixes the objects and restrictions below; changing those objects changes the question. Established identities, finite calculations, non-identification constructions and solved special cases are not being posed anew as open. Some residual tasks concern informative bounds, biological realization or external validation rather than a new theorem.

Every main definition and problem statement received a mathematical editorial review. Each entry now has an additional binding-conditions paragraph; direct corrections are separately recorded in the accompanying audit. The alternative formulations inherit the applicable conventions, but are different targets unless equivalence is proved. Where an old broad phrase conflicts with an explicit restriction below it, the explicit restriction controls.

\textbf{Literature boundary.} Thirteen selective primary-literature checks cover eighteen entries in this revision; their access levels and scope are recorded explicitly. The supplied bibliography and its earlier source-access descriptions are retained as provenance, not represented as newly verified. Neither this review nor a source's description proves the absence of later solutions. The date of this edition is 13 October 2026; there is no claim of exhaustive current-status certification for all 330 entries.

\subsection*{Binding conventions for a concrete instance}
\paragraph{Objects, quantifiers and non-vacuity.}
An instance declares measurable state, input, output and parameter spaces; a fixed candidate family; admissible initial laws, interventions and environments; a finite horizon or specified asymptotic regime; observation and sampling laws; and performance metrics, positive normalization scales and positive tolerances. These are chosen before fitting or testing, not adjusted after seeing failures. Shared parameters must genuinely be shared across the declared validation domain. Finite dimension alone does not restrict arbitrary real-number encodings: specify a finite library, bounded description length, or regularity/complexity bounds. Exclude constant or no-signal solutions through a fixed informative-output or baseline-improvement requirement where appropriate. A universally quantified result cannot be established by choosing one trivial favorable instance.

\paragraph{Probability laws and data.}
Every expectation or probability refers to the specified stochastic law, including initial uncertainty, process noise, observation noise and sampling as appropriate. A deterministic equation is not a probability model without a declared random input or initial law. ``Observed'' targets in population-level formulas mean population observables $g_j^*$, not noisy estimates asserted to be exact. Let $\widetilde g_j$ be independent validation estimates and suppose
\[
\mathbb{P}\{\lvert\widetilde g_j-g_j^*\rvert\le r_j\text{ for every }j\}\ge1-\alpha.
\]
Then an uncertainty-aware sufficient certification rule for positive target tolerances $\varepsilon_j$ is
\[
\max_j\frac{|\widehat g_j-\widetilde g_j|+r_j}{\varepsilon_j}\le1.
\]
On that simultaneous confidence event it implies every population error is at most its tolerance, by the triangle inequality. Vector/norm versions use corresponding valid simultaneous bounds. Training uncertainty must also be included when the claim averages over repeated training data. Evaluation on finitely many conditions does not establish a continuous-domain supremum without covering and regularity bounds.

\paragraph{Dynamics and well-posedness.}
For ODEs and SDEs give coefficient domains, initial laws and existence/uniqueness conditions in the intended solution sense, or keep the entire admitted solution set. For stochastic jump models specify nonnegative off-diagonal rates, diagonal normalization, rate orientation and non-explosion. For PDEs specify spatial domains, boundary/initial conditions, conservation or source terms, and a weak or classical solution class. Moving-domain concentration equations include dilution/transport terms. A stationary law is used only for a time-homogeneous augmented model under existence assumptions; a stationary approximation to a driven system must be justified. Absorbing extinction states are retained. Survival-conditioned, quasi-stationary and unconditional laws are distinct and must not be substituted for one another.

\paragraph{Identification and causal meaning.}
Two parameter values are observationally equivalent when they induce the same full distribution under all allowed observation/intervention designs, not merely the same likelihood value for one sample. A target is identified only if constant on the equivalence class. Gauge symmetries, label switching and reference conventions are quotiented or fixed. Finite-data compatible sets use a declared confidence, likelihood or discrepancy rule; an empty set signals model rejection. Bayesian credible sets are not automatically confidence sets. A prior or regularizer does not create data identification. Bounds intended to resolve a question must be informative at a predeclared width or decision threshold.

A causal contrast uses well-defined intervention versions and a stated causal model. Random assignment, or explicit consistency, overlap and exchangeability assumptions, is required for causal interpretation from observations. Interference and time-varying confounding are represented when present. Fixing a mediator by intervention differs from conditioning on its observed value. Marginal potential-outcome laws do not determine individual cross-world effects or an arbitrary joint coupling. Ethically admissible designs, rather than unconstrained interventions, define the experimental menu.

\paragraph{Metrics, units, optimization and limits.}
Specify the sampling or reference measure for densities and integrals, the log base, and dimensionless arguments of logarithms. Total variation is $\sup_A|P(A)-Q(A)|$; Wasserstein distances require the declared ground metric and appropriate moment conditions; KL divergence uses its stated absolute-continuity convention. All normed readouts are scaled consistently, and error probabilities use $0<\alpha<1$. A ratio has a positive denominator or an explicit undefined/zero convention. Thermodynamic entropy production is nonnegative in expectation under the relevant hypotheses; individual stochastic realizations need not be. An infimum is not an attained minimum without compactness/coercivity and continuity or another existence argument. Fix the order of large-population, long-time and vanishing-mutation limits. Stability claims declare perturbation space and symmetries.

\subsection*{How the provisional difficulty order was assigned}
Four editorial components are scored from 1 (lower burden) to 5 (higher burden): (a) formal or inferential difficulty, (b) access to evidence and identification burden, (c) mechanistic/computational integration, and (d) independent validation and transfer. Their unweighted total orders the entries. For selected entries in the highest bands, an explicit editorial priority breaks equal totals; all other ties are broken by original problem number. Equal totals are \emph{not} evidence of reliably different true difficulties. Scores are judgments, not measurements, completion-time estimates or consensus rankings.

The object ranked is an informative, nondegenerate resolution of the \emph{scoped} scientific target, with the stated budget and validation domain. A sharp useful impossibility result can qualify; an already known identity, vacuous bound, arbitrary favorable model or lookup table cannot. Rank 330 is merely the least difficult relative assessment in this catalogue, not a claim of an easy biological solution. Broad historical titles do not automatically outrank precisely formulated mechanistic tasks.



\end{document}
